% Section 7 (OUTLINE.md §7; 2 pages). The listings are pulled from docs/paper/code/ and their
% outputs from docs/paper/code/out/, written by docs/paper/code/run_examples.py, which ran both
% scripts with `uv run --frozen python` in the main repository at commit d58c9ef (branch
% phase-2-cr-acceptance; includes D-34, b9c66c0) on macOS-26.6.2-arm64 (versions in
% code/out/env.json). The main working tree had uncommitted edits in docs/phase2/DESIGN.md,
% src/rai_bvar/validation/joint.py and test files only (listed in env.json); no code path of the
% examples imports rai_bvar.validation (grep at d58c9ef: the only production import is inside the
% prior mode of structural/ordering.py, which the examples do not call). Both outputs are identical, line for
% line, to those recorded at 62d15f6 in dossier 02 §13 (ex_b without its two wall-clock lines,
% which were removed from the script). Re-run on the release commit (TODO 20).
\section{Using \pkg{rai-bvar}: a walk-through}\label{sec:walkthrough}

This section shows the user interface on two short scripts.
Both use the same simulated data, a stationary three-variable quarterly VAR(1) for output growth,
inflation and an interest rate, stored as raw levels with FRED-style transformation codes, so they
run without network access.
\cref{lst:Lst1} uses the analytic engine and \cref{lst:Lst2} the MCMC engine; each is
followed by its standard output, verbatim.
The scripts and their outputs are in the paper's repository; the analytic output is deterministic
given the seeds, and the MCMC output is bitwise reproducible only on the same platform and package
versions (\cref{sec:architecture}).
% [V] docs/paper/code/{ex_a,ex_b}.py, out/{ex_a,ex_b}.txt, out/env.json (d58c9ef, macOS arm64); identical to dossier 02 §13 outputs at 62d15f6
\todo{Length: with both listings and their outputs this section runs to about four pages in the
assembled build against the two budgeted in OUTLINE \S7. Options: move the two output blocks to \cref{app:algorithms},
or typeset each script and its output side by side.}

\subsection{Data}\label{sec:walk-data}

\code{VARData.from_frame(raw, tcodes={...})} takes a data frame of raw levels indexed by
dates or periods.
Its column order is the recursive ordering used by every triangular model.
Gaps in the index, duplicate columns, infinite values, logs of non-positive values and missing
values inside the sample raise errors; leading and trailing incomplete periods are trimmed, with a
warning except for the leading periods that the transformations themselves consume.
% [V] tests/unit/test_dataset.py::test_interior_missing_value_raises; dossier 02 §9.1; src/rai_bvar/data/dataset.py l.215-237 (_trim warns on leading periods only when first > the transforms' history length) at d58c9ef
\code{upto(origin)}, \code{tail(n)}, \code{select(variables)} and \code{realized(origin, h)} give
the data available at an origin, rolling windows, reorderings and future outcomes on the
transformed scale.
The transformation codes 1--7 are those of FRED-MD and FRED-QD \citep{mccracken2016fredmd}, each
with a pathwise inverse, and they set the own-first-lag prior mean (1 for levels and log levels,
0 for differences; \cref{sec:framework}).
For real data, \code{rb.datasets.load("fred-qd", vintage="2026-09")} downloads a pinned vintage
of FRED-QD \citep{mccracken2021fredqd} on demand, checks its SHA-256 checksum and caches it; the
files are not bundled with the package.
Its \code{variables} argument takes a named set or a list of series (the order given is the
ordering), and \code{tcodes} takes \code{"rai"}, \code{"official"} or a mapping.
The named variable sets (\code{"small"}, \code{"medium"}, \code{"large"}) contain only series
published by U.S.\ federal agencies, and interior publication gaps are interpolated (linearly, in
logs for logged series) with a warning that lists every filled period (\code{fill\_gaps=True}).
% [V] src/rai_bvar/datasets/fred.py module docstring and load() (l.169-179) at d58c9ef; dossier 02 §9.5
\todo{Q11/TODO 21: confirm the official FRED-QD 2026-09 transformation codes against the
\code{"rai"} codes before the paper states them.}

\subsection{An analytic BVAR, a forecast and a comparison}\label{sec:walk-analytic}

\lstinputlisting[language=Python, firstline=6, firstnumber=6, numbers=left, numberstyle=\tiny,
  basicstyle=\scriptsize\ttfamily, columns=fullflexible, keepspaces=true, breaklines=true,
  showstringspaces=false, upquote=true,
  caption={\code{ex\_a.py}: data, an analytic BVAR with SOC/SUR dummies and GLP empirical Bayes,
  a forecast, and a rolling-window comparison with its verdicts.},
  label=lst:Lst1]{code/ex_a.py}
\lstinputlisting[basicstyle=\scriptsize\ttfamily, columns=fullflexible, keepspaces=true,
  breaklines=true, upquote=true, title={Output of \cref{lst:Lst1} (stdout, verbatim; nothing
  was written to stderr)}]{code/out/ex_a.txt}

\cref{lst:Lst1} runs on the analytic engine.
The log difference (code 5) consumes the first period (239 of 240 raw rows remain), and the
own-first-lag prior means follow the codes: 0, 0 and 1.
% [V] docs/paper/code/out/ex_a.txt lines 1-2
The model of line 34 adds SOC and SUR dummy observations to the Minnesota prior and leaves the
volatility at its default, so the composition resolves to the conjugate NIW model on the analytic engine, and
empirical Bayes in the hierarchy of \citet{giannone2015prior} sets the tightness (printed as
\code{lambda}, $\kappa$ in \cref{sec:framework}) and the SOC and SUR weights.
% [V] out/ex_a.txt lines 3-7 (kappa 0.268, soc 1.120, sur 1.177)
The forecast is summarised on the transformed scale and, through the pathwise inverse, on the
level scale.

The comparison runs 39 rolling-window origins at $h = 1$ and 36 at $h = 4$.
% [V] out/ex_a.txt table, column origins
The OLS VAR(2) is ``better'' than the random walk along the \code{model} axis at $h = 1$: its
$\Delta$LPD is $+14.45$ with an MCSE of 0.10, the one-sided DM--HLN $p$-value is 0.014
on a rolling window, and the MCS keeps the VAR and excludes the random walk.
% [V] out/ex_a.txt, reasons line "OLS-VAR(2) vs RW h=1"
The table prints the two-sided $p$-value, 0.029, for the same pair.
% [V] out/ex_a.txt table row "OLS-VAR(2) 1"
The comment in the listing says ``DM/AG''; the test implemented is the Diebold--Mariano test with
the Harvey--Leybourne--Newbold correction applied to log-score losses (\cref{sec:eval-tests}).
% risk 10: the listing is kept verbatim
The BVAR has the larger $\Delta$LPD ($+15.37$ at $h = 1$), yet its verdict against the random walk
is ``not comparable'': the two differ along seven axes, and a verdict needs exactly one (R2).
% [V] out/ex_a.txt table row "BVAR(2) 1" and the verdict lines
A comparison that credits the prior would pair the BVAR with an arm that differs from it only in
the prior.

\subsection{The same model family on the MCMC engine}\label{sec:walk-mcmc}

\lstinputlisting[language=Python, firstline=37, lastline=75, firstnumber=37, numbers=left,
  numberstyle=\tiny, basicstyle=\scriptsize\ttfamily, columns=fullflexible, keepspaces=true,
  breaklines=true, showstringspaces=false, upquote=true,
  caption={\code{ex\_b.py} (lines 37--75; lines 8--36 hold the imports, the \pkg{pandas} display
  options and the data block of \cref{lst:Lst1}): stochastic volatility on the MCMC engine, summary, forecast LPD,
  impulse responses, and a volatility-axis comparison.},
  label=lst:Lst2]{code/ex_b.py}
\lstinputlisting[basicstyle=\scriptsize\ttfamily, columns=fullflexible, keepspaces=true,
  breaklines=true, upquote=true, title={Output of \cref{lst:Lst2} (stdout, verbatim; MCMC
  numbers are platform-specific)}]{code/out/ex_b.txt}

With the volatility set to \code{SV()}, the default mean, \code{Linear(prior=Minnesota())} (without
the dummy observations of \cref{lst:Lst1}), moves to the MCMC engine, in structural form with
the Mode~I prior; there is no engine argument to set.
% [V] out/ex_b.txt line 2; tests/unit/test_api_resolution.py::test_resolution_table
The Minnesota hyperparameters are an empirical-Bayes plug-in from the conjugate twin on the same
window (\code{eb\_method: conjugate\_twin}), so the mean prior is the one every volatility arm on
that window receives; for the SV model this is a plug-in value, not a type-II maximum-likelihood
estimate.
% [V] out/ex_b.txt line 3; tests/analytic/test_eb_mcmc.py::test_eb_plug_in_equals_the_phase1_eb_mode_at_cross_1
The summary reports $\hat R$ and bulk ESS for every primary quantity.
On this homoskedastic DGP the SV parameters mix slowly: $\phi_i$ and $\sigma_i$ have bulk ESS
between 110 and 362 from 4000 kept draws and $\hat R$ up to 1.042, so the fit is flagged as not
converged; the library's warning (on stderr, not shown) reports that 7 of the 16 primary
quantities fail, which by the table are the six SV parameters $\phi_i$, $\sigma_i$ and the
log-likelihood ($\hat R$ 1.031, bulk ESS 179).
% [V] out/ex_b.txt summary table and "converged: False" (MCMCSettings(chains=4, draws=1000): 4 x 1000 kept draws); out/ex_b.stderr.txt ("7 of 16 primary quantities fail", naming the five worst; warning text not reproduced; a second warning names other packages, TODO 16)
The forecast container and its log-density interface are those of the analytic case (R3); the
\code{engine} attribute records which engine produced it, and the MCMC forecast also carries the
ordering, sampler metadata and accelerator variables (\cref{sec:eval-contract}).
% [V] out/ex_b.txt lines "ForecastResult(...) | engine attr: mcmc" and the two LPD lines
With the interest rate ordered last, a one-s.d.\ rate shock has no impact effect on output growth
or inflation in the recursive identification.
% [V] out/ex_b.txt IRF table, horizon 0 (0.0000, 0.0)

The final comparison is along the volatility axis only: both arms are structural Mode~I models with
the same mean prior, and their axis labels differ only in \code{vol}.
% [V] out/ex_b.txt axes lines
Over 11 rolling-window origins, SV has a $\Delta$LPD of $-0.11$ with an MCSE of 0.10,
too imprecise for any verdict, and the one-sided DM $p$-value is 0.421.
% [V] out/ex_b.txt table row "Linear x SV 1" and reasons line
The SV arm converged at none of the 11 origins, so the convergence gate forces
``inconclusive'' whatever the other evidence says.
% [V] out/ex_b.txt reasons line "converged at only 0% of origins"
This is the intended behaviour: the gate refuses to certify a comparison with an unconverged fit.
A longer chain is the first remedy to try; \cref{sec:cr} shows how the gate decided a criterion
of the replication and the fit-level diagnostic that set the MCMC budget before the acceptance run.

\subsection{Inspecting fits and failures}\label{sec:walk-inspect}

An MCMC fit exposes its Gibbs plan: \code{fit.plan.describe()} prints the order of one sweep with
each step's block and write set (\cref{fig:FigC1} in \cref{app:algorithms}).
% [V] src/rai_bvar/engines/mcmc/plan.py l.91-103 at d58c9ef
It also exposes the diagnostics table, \code{summary()}, the \code{converged} flag, the list of
\code{warnings} and the timing of each stage.
A numerical failure inside a kernel reaches the user as \code{MCMCNumericalError}, which names the
block, the equation, the chain and the sweep and carries the diagnostics of the deterministic replay
(\cref{sec:arch-mcmc}; \cref{sec:robustness} describes the failures found this way).
% [V] src/rai_bvar/engines/mcmc/protocols.py l.165-189 at d58c9ef
\todo{OUTLINE \S7.4: a volatility-axis grid, \code{rb.grid(rb.BVAR, vol=[rb.Constant(), rb.SV(),
rb.SV(outliers=True)])} with \code{verdict(axis="vol")}, is [TC]. Add it as a third script under
\code{docs/paper/code/} and to \code{run\_examples.py} only if it is run on the release commit.}
